
\documentclass{article}
\usepackage{colortbl}
\usepackage{makecell}
\usepackage{multirow}
\usepackage{supertabular}

\begin{document}

\newcounter{utterance}

\centering \large Interaction Transcript for game `cladder', experiment `full\_v1.5\_default', episode 1876 with Qwen3.5{-}4B.
\vspace{24pt}

{ \footnotesize  \setcounter{utterance}{1}
\setlength{\tabcolsep}{0pt}
\begin{supertabular}{c@{$\;$}|p{.15\linewidth}@{}p{.15\linewidth}p{.15\linewidth}p{.15\linewidth}p{.15\linewidth}p{.15\linewidth}}
    \# & \multicolumn{2}{c}{Player} && \multicolumn{2}{c}{Game Master} \\
    \hline

    \theutterance \stepcounter{utterance}  
    & & & \multicolumn{4}{p{0.6\linewidth}}{
        \cellcolor[rgb]{0.9,0.9,0.9}{
            \makecell[{{p{\linewidth}}}]{
                \texttt{\tiny{[P1$\langle$GM]}}
                \texttt{You are an expert in causal inference. The following question is not a typical commonsense query, but rather a meticulously designed question created by a professor specializing in causal inference, intended to assess the students' mastery of the course content.} \\
\texttt{Imagine a self{-}contained, hypothetical world with only the following conditions, and without any unmentioned factors or causal relationships: Medication has a direct effect on blood pressure and heart condition. Blood pressure has a direct effect on heart condition. For people not taking any medication, the probability of healthy heart is 64\%. For people taking medication, the probability of healthy heart is 21\%. For people taking medication, would it be more likely to see healthy heart if the person had taken no medication?} \\
\texttt{Start your answer with 'Yes' or 'No', followed by additional reasoning or evidence to support your explanation.} \\
            }
        }
    }
     \\ \\

    \theutterance \stepcounter{utterance}  
    & \multicolumn{4}{p{0.6\linewidth}}{
        \cellcolor[rgb]{0.561,0.737,0.561}{
            \makecell[{{p{\linewidth}}}]{
                \texttt{\tiny{[P1$\rangle$GM]}}
                \texttt{No} \\
\\ 
\texttt{**Reasoning:**} \\
\\ 
\texttt{This question is a classic trap designed to test the distinction between **marginal probabilities** (observed frequencies) and **counterfactual probabilities** (what would have happened under a different intervention).} \\
\\ 
\texttt{1.  **Analyze the Causal Structure**:} \\
\texttt{    *   Let \$M\$ be the treatment (Medication).} \\
\texttt{    *   Let \$BP\$ be the mediator (Blood Pressure).} \\
\texttt{    *   Let \$H\$ be the outcome (Healthy Heart).} \\
\texttt{    *   The causal chain is \$M \textbackslash{}rightarrow BP \textbackslash{}rightarrow H\$.} \\
\texttt{    *   Crucially, the prompt states: "Medication has a direct effect on blood pressure and heart condition." This implies \$M\$ also has a direct arrow to \$H\$ (\$M \textbackslash{}rightarrow H\$), in addition to the indirect path through \$BP\$.} \\
\\ 
\texttt{2.  **Analyze the Given Probabilities**:} \\
\texttt{    *   \$P(H=1 | M=0) = 0.64\$ (Probability of healthy heart for those *not* taking medication).} \\
\texttt{    *   \$P(H=1 | M=1) = 0.21\$ (Probability of healthy heart for those *taking* medication).} \\
\\ 
\texttt{3.  **Evaluate the Counterfactual Question**:} \\
\texttt{    *   The question asks: "For people taking medication (\$M=1\$), would it be more likely to see a healthy heart if the person had taken no medication (\$M=0\$)?"} \\
\texttt{    *   Mathematically, this asks if \$P(H=1 | do(M=1), M=0) > P(H=1 | do(M=0))\$.} \\
\texttt{    *   In causal inference, the term \$P(H=1 | do(M=1), M=0)\$ represents the counterfactual: "What is the probability of a healthy heart for a person who is actually taking medication, *if we had intervened* to make them take no medication?"} \\
\texttt{    *   By the definition of the counterfactual operator \$do(M=0)\$, this probability collapses to the observed probability for the group that takes no medication. Therefore, \$P(H=1 | do(M=1), M=0) = P(H=1 | M=0)\$.} \\
\\ 
\texttt{4.  **Compare the Values**:} \\
\texttt{    *   The probability of a healthy heart if the person had taken no medication is **64\%**.} \\
\texttt{    *   The probability of a healthy heart if the person is taking medication is **21\%**.} \\
\texttt{    *   Since \$64\textbackslash{}\% > 21\textbackslash{}\%\$, the outcome is indeed *more likely* under the counterfactual scenario (\$M=0\$) than under the actual scenario (\$M=1\$).} \\
\\ 
\texttt{5.  **Address the "Yes/No" Ambiguity**:} \\
\texttt{    *   The question asks:} \\
            }
        }
    }
    & & \\ \\

    \theutterance \stepcounter{utterance}  
    & & & \multicolumn{2}{p{0.3\linewidth}}{
        \cellcolor[rgb]{0.9,0.9,0.9}{
            \makecell[{{p{\linewidth}}}]{
                \texttt{\tiny{[GM$|$GM]}}
                \texttt{No} \\
            }
        }
    }
    & & \\ \\

    \theutterance \stepcounter{utterance}  
    & & & \multicolumn{2}{p{0.3\linewidth}}{
        \cellcolor[rgb]{0.9,0.9,0.9}{
            \makecell[{{p{\linewidth}}}]{
                \texttt{\tiny{[GM$|$GM]}}
                \texttt{yes} \\
            }
        }
    }
    & & \\ \\

    \theutterance \stepcounter{utterance}  
    & & & \multicolumn{2}{p{0.3\linewidth}}{
        \cellcolor[rgb]{0.9,0.9,0.9}{
            \makecell[{{p{\linewidth}}}]{
                \texttt{\tiny{[GM$|$GM]}}
                \texttt{game\_result = LOSE} \\
            }
        }
    }
    & & \\ \\

\end{supertabular}
}

\end{document}
